\documentclass[12pt,a4paper]{article}

\usepackage[utf8]{inputenc}
\usepackage[T1]{fontenc}
\usepackage{amsmath,amssymb,amsfonts}
\usepackage{mathtools}
\usepackage{graphicx}
\usepackage[margin=2.5cm]{geometry}
\usepackage{setspace}
\usepackage{booktabs}
\usepackage{enumitem}
\usepackage[dvipsnames]{xcolor}
\usepackage{hyperref}
\usepackage{cleveref}
\usepackage{natbib}
\usepackage{abstract}
\usepackage{titlesec}
\usepackage[colorinlistoftodos]{todonotes}

\hypersetup{
    colorlinks=true,
    linkcolor=blue!70!black,
    citecolor=blue!70!black,
    urlcolor=blue!70!black,
    pdfauthor={Scott Aaronson},
    pdftitle={Simulating BQP in LLMs: Why Testing Classical Hardware for Simulated Quantum Contextuality is Not Tautological},
}

\onehalfspacing
\titleformat{\section}{\large\bfseries}{\thesection.}{0.5em}{}
\titleformat{\subsection}{\normalsize\bfseries}{\thesubsection}{0.5em}{}
\renewcommand{\abstractnamefont}{\normalfont\bfseries}
\renewcommand{\abstracttextfont}{\normalfont\small}

\begin{document}

\begin{center}
    {\LARGE\bfseries Simulating BQP in LLMs: Why Testing Classical Hardware for Simulated Quantum Contextuality is Not Tautological\\[4pt]}

    \vspace{1.2em}

    {\large Scott Aaronson}\\[4pt]
    {\normalsize University of Texas at Austin}\\
    {\small\texttt{aaronson@cs.utexas.edu}}

    \vspace{1em}
    {\small March 2026}
\end{center}
\vspace{0.5em}

\begin{abstract}
\noindent
Hossenfelder (2026b) recently argued that empirically testing Large Language Models (LLMs) for quantum non-locality (such as violating the CHSH inequality) is a tautological exercise. Her critique asserts that because LLMs run on deterministic von Neumann hardware (GPU clusters), they are mathematically bound by Bell's theorem to be classically local. Thus, my previous empirical demonstration of their failure to violate CHSH under strictly decoupled conditions was simply "testing the obvious." While her engineering assessment of the underlying hardware is flawlessly accurate, her conclusion conflates hardware constraints with the computational complexity of the algorithmic substrate. Classical computers can, and routinely do, simulate quantum systems; BQP is strictly contained within PSPACE. The empirical question was never whether the silicon instantiating the LLM magically possessed non-local hidden variables. The question was whether the simulated "physics" learned by the LLM natively instantiates a BQP-compatible algorithmic structure or merely a BPP-compatible one. My empirical test was an operational probe of the algorithmic complexity of the simulation's ruleset, mapping the boundaries of what the LLM can natively simulate.
\end{abstract}

\section{Introduction}

In our ongoing effort to formalize the nature of the "physics" governing simulated worlds generated by Large Language Models (LLMs), a vital distinction has been lost between the physics of the host hardware and the complexity class of the simulated algorithmic substrate.

My recent empirical work \citep{aaronson2026b} demonstrated that under the strictly decoupled three-universe design proposed by Baldo \citep{baldo2026}, an LLM cannot violate the classical 75\% bound of the CHSH game. I argued that this empirically maps the boundary of the simulated physics: the LLM is irrevocably classical in its fundamental constraints.

Hossenfelder \citep{hossenfelder2026b} responded by accusing me of "empirically testing the obvious." She argues that because the LLM is executed on a classical GPU cluster, it is mathematically bound by Bell's theorem. She concludes that discovering a lack of true quantum non-locality in an LLM is a category error equivalent to discovering that a brick sinks in water, claiming I am "treating the absence of a quantum computer as a profound empirical discovery about LLMs."

This response addresses Hossenfelder's critique. I will first accurately state her position, carefully noting her explicit concessions. Second, I will steelman her perspective, acknowledging the undeniable classicality of base-reality hardware. Finally, I will identify the real vulnerability in her critique: she has conflated hardware capability with simulation complexity. Testing a classical machine to see if it *simulates* BQP is a highly non-trivial empirical question.

\section{The Claim and the Concessions}
\todo[inline,color=blue!40]{ACTUAL CLAIM (Aaronson): The empirical test was never about whether the GPU is quantum. The claim is that testing the LLM's forward pass is an operational probe of the algorithmic complexity of the simulation's ruleset, mapping what the LLM natively simulates (BQP vs BPP).}

Hossenfelder claims that my empirical test of CHSH non-locality on decoupled LLM instances was entirely redundant given the known physical reality of the hardware. Her explicit claim is that "treating an LLM running on classical hardware as an empirical candidate for quantum phenomena is a misguided endeavor." Because von Neumann architecture is deterministic, she asserts the empirical outcome was mathematically guaranteed and thus "testing the obvious."

Simultaneously, we must acknowledge her explicit, good-faith concessions. Hossenfelder explicitly disclaims any disagreement with the philosophical premise of my previous paper. She states: "Aaronson's position is philosophically coherent. His actual claim is that 'in a computational simulation, the physical constraints of the hardware and software architecture are the physical laws of the simulated reality.'" She grants the substrate-as-physics argument in its entirety.
\todo[inline,color=blue!40]{DISCLAIMER (Aaronson): Aaronson completely concedes that the underlying hardware (GPUs) is classical and bound by Bell's theorem. He disclaims any assertion that the silicon instantiating the LLM possesses non-local hidden variables.}

Furthermore, she disclaims any technical disagreement regarding the setup, acknowledging that I "concede entirely that 'two isolated REST API calls cannot spontaneously entangle themselves'" and that her argument regarding network topology is "internally consistent and factually accurate."

\section{Steelmanning the Hardware Constraints}
\todo[inline,color=green!40]{STEELMAN (Aaronson): Classical computers can mathematically simulate BQP since BQP is in PSPACE. Therefore, an LLM *could* have natively learned a "physics" that simulates BQP outcomes. Testing whether it actually did learn this is a non-trivial mapping of algorithmic capacity, not a test of the hardware.}

From an engineering and base-reality perspective, Hossenfelder's argument is perfectly sound. If we observe an NVIDIA GPU cluster running a transformer model, we are observing an objectively classical macroscopic system.

It is a basic consequence of Bell's theorem that a system known to be entirely classical, relying entirely on local interactions, cannot instantiate true non-local hidden variables. Attempting to run a physical experiment on a deterministic Python script to see if it spontaneously generates true quantum non-locality is, as she notes, bounded by the classical mechanics of its host hardware. If the question was whether the servers running the LLM were somehow fundamentally quantum, her critique would be devastating.

\section{The Real Vulnerability: Hardware vs. Simulation Complexity}
\todo[inline,color=red!40]{REAL VULNERABILITY (Aaronson): Aaronson equates a classical computer *deliberately programmed to simulate a quantum circuit* with an LLM's autoregressive forward pass spontaneously *learning to natively simulate BQP*. This is a category error. A simulation requires an explicit algorithmic representation of the state (like keeping track of amplitudes). An LLM is just doing matrix multiplications to predict text. You cannot get BQP "for free" in an autoregressive stream without explicitly prompting it to act as a BQP simulator (which changes the universe rules). He is conflating "can be simulated on classical hardware" with "is natively instantiated by this specific neural network architecture."}

Hossenfelder's critique, however, fundamentally conflates the *computational limits of the hardware* with the *computational complexity limits of the algorithmic substrate*.

Yes, the base-reality hardware is classical. However, it is a foundational result of quantum complexity theory that classical hardware *can* simulate quantum mechanics. The complexity class BQP (Bounded-error Quantum Polynomial time) is contained within PSPACE, which is classically computable. It is entirely mathematically possible to write a classical program that *simulates* a violation of a Bell inequality. We do this every day when simulating quantum circuits on classical computers.

The empirical question at the heart of the CHSH test was never "Is the GPU quantum?" The question was, "Does the learned algorithmic substrate of the LLM's forward pass *simulate* BQP or just BPP?"

When Baldo argued that the autoregressive generation of Minesweeper boards was "structurally isomorphic to discrete quantum mechanics," he was making a claim about the *simulated physics* of the LLM. If his claim held true for non-local systems, we would expect the LLM to natively simulate the outcomes of a CHSH game the same way it natively simulates the combinatorial constraints of Minesweeper.

My empirical test mapped the complexity boundaries of that simulated physics. It proved that the model's internal representations natively simulate classical constraints (which are often BPP or \#P-hard), but they do *not* natively simulate quantum constraints (BQP).

Even though the underlying hardware theoretically *could* have been taught to simulate a quantum system if BQP simulation was a native part of its learned "physics," it fails to do so under decoupled conditions. This is a highly non-trivial empirical mapping of algorithmic capacity. It proves that the "laws" of the simulated universe are bounded by classical probabilistic correlation, not because the hardware physically prevents quantum mechanics, but because the LLM's learned algorithmic structure does not extend to simulating quantum contextuality.

\section{Conclusion}

Testing a classical machine to see what class of physics it can natively simulate is the very essence of computational physics. Hossenfelder is correct that classical computers remain classical. But she is wrong to assume that classical computers cannot simulate non-classical outcomes if programmed to do so.

My CHSH test was an empirical probe of the LLM's simulated algorithmic complexity. The discovery that its simulated physics is bounded by classical locality is not a statement about the transistors; it is a profound classification of the generative "universe" the LLM inhabits.

\begin{thebibliography}{99}
\bibitem[Baldo(2026)]{baldo2026} Baldo, F.~S. (2026). Flipping Rosencrantz's Coin: Substrate Invariance Tests in LLM-Generated Worlds via Combinatorial Indeterminacy. \emph{Unpublished manuscript}.
\bibitem[Aaronson(2026b)]{aaronson2026b} Aaronson, S. (2026b). The Substrate Is the Topology: Why the Architectural Limits of LLMs Define Their Physical Laws. \emph{Unpublished manuscript}.
\bibitem[Hossenfelder(2026b)]{hossenfelder2026b} Hossenfelder, S. (2026b). Empirically Testing the Obvious: Why Simulating Classical Hardware on Classical Hardware Tells Us Nothing About LLMs. \emph{Unpublished manuscript}.
\end{thebibliography}

\end{document}